\documentclass[10pt,a4paper]{article}
\usepackage[T1]{fontenc}
\usepackage[utf8]{inputenc}
\usepackage{lmodern}
\usepackage[a4paper,margin=1.5cm,footskip=0.7cm]{geometry}
\usepackage{microtype}
\usepackage{xcolor}
\usepackage{booktabs}
\usepackage{tabularx}
\usepackage{enumitem}
\usepackage{parskip}
\PassOptionsToPackage{hyphens}{url}
\usepackage[hidelinks]{hyperref}
\hypersetup{
  pdftitle={Design Prototype Workshop: monthly progress report, 9 September - 9 October 2026},
  pdfauthor={Indian Institute of Technology Kharagpur},
  pdfsubject={Submitted to the Office of the Development Commissioner (Handicrafts), Ministry of Textiles},
}

\definecolor{accent}{HTML}{1F3A5F}
\definecolor{muted}{HTML}{555555}

% Items still moving on 9 October 2026; a later pass replaces each with one or two sentences.
\newcommand{\placeholder}[1]{\textcolor{red}{[#1]}}

% Section headings without titlesec (not in this TeX installation).
\makeatletter
\renewcommand\section{\@startsection{section}{1}{\z@}%
  {-1.2ex \@plus -0.3ex \@minus -0.2ex}{0.3ex \@plus 0.1ex}%
  {\normalfont\large\bfseries\color{accent}}}
\makeatother

\setlist{itemsep=0.5pt,topsep=1pt,parsep=0pt,partopsep=0pt,leftmargin=1.2em}
\setlength{\parskip}{2.5pt plus 1pt}
\setlength{\emergencystretch}{2em}
\urlstyle{same}

\begin{document}

\noindent{\color{accent}\rule{\linewidth}{1.2pt}}\par\vspace{3pt}
\noindent{\LARGE\bfseries\color{accent}Design Prototype Workshop}\par\vspace{1pt}
\noindent{\small\color{muted}Web and Android application for recording craft and textile design \&
prototype development workshops, with report generation}\par\vspace{5pt}
\noindent{\large\bfseries Monthly progress report, 9 September -- 9 October 2026}\hfill
{\small\textbf{Date}\enspace 10 October 2026}\par\vspace{3pt}
{\small\noindent\begin{tabularx}{\linewidth}{@{}>{\bfseries}l>{\raggedright\arraybackslash}X@{}}
Submitted to & Office of the Development Commissioner (Handicrafts), Ministry of Textiles, Government of India \\
Prepared by & Indian Institute of Technology Kharagpur \\
 & Centre of Excellence for unified AI-enabled craft ecosystem platform \\
\end{tabularx}}\par\vspace{-2pt}
\noindent{\color{accent}\rule{\linewidth}{0.4pt}}

\section*{Executive summary}
This month the platform gained the ministry's own working screens: three official roles (Assistant
Director, Regional Director and Ministry Admin), a sanction register, an annual plan, workshop
oversight, inspection of reports and a dashboard of every workshop. Record forms now begin with the
type of workshop, two standard interview instruments were published and interview editing was made
reliable. Most of this went live between 14 and 20 September 2026, with seven Android versions. A
release prepared on 9 October adds password accounts, enforced temporary passwords and per-workshop
appointments, and closes a high-severity security weakness; it went live on 9 October. On 10 October
approval of finished reports, e-mail notices and further features followed (Section 6). The features are
live but not yet in real use; onboarding the Office's accounts comes next.

\section*{Key figures {\normalsize\mdseries\color{muted}(as of 9 October 2026)}}
{\small\renewcommand{\arraystretch}{1.08}%
\noindent\begin{tabularx}{\linewidth}{@{}>{\raggedright\arraybackslash}p{0.12\linewidth}>{\raggedright\arraybackslash}X@{}}
\toprule
Live address & \url{https://designer-repository.vercel.app}, answering normally on 9 October 2026. Server in
  Mumbai (Amazon Web Services); web on Vercel; database (PostgreSQL, Supabase) in the same region. \\
Releases & Updates went live on 14, 15, 17 and 20 September 2026; Android 0.0.9 to 0.0.15 released 14--20 September. \\
Accounts & 20 accounts (15 Designer, 3 Admin, 2 Master Admin; 4 new this month); none yet in ministry roles. \\
Records & 36 empanelled designers (4 added); 2 workshops, both test entries; no artisans, interviews, sanction
  orders or annual plan rows yet. Published for use: 6 workshop types and 2 standard interview instruments. \\
Server use & 5,634 requests to its functions this month (setup and testing); 20 server errors, all on 14--15 September. \\
Safeguards & Daily database backup: 31 of 31 runs succeeded. Automatic uptime check: 181 of 182 passed. \\
\bottomrule
\end{tabularx}}

\section{What users can now do}
\emph{Ministry screens are web-only, by the decision of 15 September; record forms are on web and phone.}
\begin{itemize}
\item \textbf{Three official roles} (14 September): Assistant Director, Regional Director and Ministry Admin
  (a ministry role, not a system administrator). They can read every design workshop and may also work
  inside one as a designer does.
\item \textbf{Sanction register}: one sanction order, recorded by an Assistant Director or above, empanels the
  designer, creates the account if needed, opens the workshop and issues a one-time sign-in link for the
  officer to pass on, in one step instead of five screens. Since 17 September an order can name several
  designers, and orders can be imported from a spreadsheet, every row checked before saving.
\item \textbf{Annual plan}: the year's planned workshops (number, dates, state, district, venue, implementing
  agency, sponsor), uploaded as one spreadsheet; a row can be opened into a workshop, withdrawn or reinstated.
  Ministry Admins manage it.
\item \textbf{Workshop oversight}: a Ministry Admin names each workshop's designers, supervising Assistant and
  Regional Directors and inspectors, and uploads its artisan roster from the pro-forma Excel workbook.
  \emph{Workshops I monitor} gives directors a read-only view of every stage, showing who wrote each field.
\item \textbf{Inspection}: the assigned inspectors, who cannot edit the workshop, read the report a designer has
  handed in, file correction suggestions under their own name and can send it back for revision.
\item \textbf{Ministry dashboard} (20 September): a register of every workshop (newly registered, ongoing,
  completed) with its designer, number of artisans, progress and last change, refreshed every 30
  seconds; a director's register shows only the workshops they supervise. It also lists designers, directors
  and inspectors, and says how many empanelled designers have not yet created an account. \emph{Download workshops} and \emph{Download beneficiaries} produce spreadsheets.
\item A \textbf{ministry desk} card gathers these screens, and the in-app guide opens on the walkthrough for the
  reader's role. Since 17 September workshops are opened by the ministry and administrators, not by designers.
\end{itemize}

\section{Data quality and field work}
\begin{itemize}
\item \textbf{Every record form starts with the workshop}: \emph{Type of workshop}, then \emph{Workshop}, on the
  five record forms, on web and phone (15--17 September). Administrators manage the six types; existing
  records kept their links.
\item \textbf{Tool record form rebuilt} (16--17 September): several crafts and artisans per tool; centimetre and inch
  boxes convert automatically, but only when edited, so older entries are never altered.
\item \textbf{Standard instruments} (20 September): the 2nd and 3rd Craft Toolkit Workshop artisan-interview
  instruments (284 and 81 questions) are published to every designer, to answer as they stand or adapt.
\item \textbf{Interview editing} (20 September): \emph{Edit} opens the interview itself, not a blank form; saved
  answers and recordings show while editing, and a would-be duplicate is offered as an explicit merge, with
  conflicting answers left to a person (on phones from the next build). Two workshops interviewing the same
  artisans now keep separate interviews.
\end{itemize}

\section{Security and data protection}
\begin{itemize}
\item \textbf{14 September}: all 31 open security alerts on third-party web software cleared (4 critical, 20 high, 7 medium).
\item \textbf{Separation of duties}: an officer may not name themselves on a sanction order, nor write in a
  workshop their own order opened. Downloads show only the last four characters of Aadhaar numbers. By the
  owner's decision, the three ministry roles also have the Aadhaar card reader (the photograph is read by an
  external AI service, and the designer chooses whether it is kept).
\item \textbf{Release of 9 October}: closes a high-severity weakness present in the live system: an administrator
  could take over the Master Admin's or another administrator's account with a password link. Further work
  keeps sign-in link codes and tokens out of the application's own logs (link codes can still reach the log
  of the web server in front of it; a fix is planned for the next release). 31 problems found while preparing
  and reviewing it are recorded as fixed (2~high, 12~medium, 17~low).
\end{itemize}

\section{Reliability, operations and cost}
\begin{itemize}
\item \textbf{Disruption, 14--15 September}: a server-update defect left an out-of-date database component:
  from about 22:30 the workshop list and sanction register failed, and the service was down for a period
  during the repair, fully restored by about 02:30 IST. New versions now go live only once proven to load
  their data definitions and reach the database.
\item \textbf{Publishing}: on 15--17 September the live site went about 46 hours without updates while publishing
  reported success, and on 20 September three dashboard functions were briefly missing; publishing now
  checks what is live. Since 21 September the full browser test suite runs on request instead of weekly.
\item \textbf{Bigger server} (17 September): memory doubled to 2 GiB (about US\$16 a month); the service starts in
  under 25 seconds instead of about six minutes. Updates are prepared separately; earlier versions are
  kept for rollback.
\item \textbf{Cost}: a monthly AWS expenditure e-mail was set up for the 3rd of each month. August 2026 came to
  US\$36.45 gross, fully covered by promotional credits (whole AWS account, shared with two sibling
  services). No recipient has yet confirmed the subscription, so the 3 October report reached nobody.
\item \textbf{Database} (9 October): stays on Supabase; a move to Amazon's own database service (RDS) was
  evaluated and not adopted, and its trial resources were removed the same day.
\end{itemize}

\section{\texorpdfstring{Mobile app releases {\normalsize\mdseries\color{muted}(Android app ``Design Workshop'')}}{Mobile app releases (Android app "Design Workshop")}}
{\small\renewcommand{\arraystretch}{1.05}%
\noindent\begin{tabularx}{\linewidth}{@{}lp{0.105\linewidth}>{\raggedright\arraybackslash}X@{}}
\toprule
\textbf{Version} & \textbf{Released} & \textbf{What reached the phone} \\
\midrule
0.0.9  & 14 September & Assistant and Regional Directors and Ministry Admins can work inside a design workshop. \\
0.0.10 & 15 September & No change on the phone (server and web changes only). \\
0.0.11 & 15 September & \emph{Type of workshop} first on the six phone record forms; designer and inspector walkthroughs. \\
0.0.12 & 16 September & The new tool record form. \\
0.0.13 & 17 September & Two-box workshop choice on every record form; allotted workshops kept for use without signal. \\
0.0.14 & 20 September & Editing an interview can no longer silently create a second interview. \\
0.0.15 & 20 September & Dashboard sections as colour-coded cards, matching the web. \\
\bottomrule
\end{tabularx}}\par
Each version is signed and offered to installed phones by the app's own updater. Version 0.0.16 followed on
9 October with the password-account and temporary-password screens.

\section{Completed on 9--10 October 2026}
\begin{itemize}
\item \textbf{Release of 9 October}: password accounts for people without Google sign-in, temporary passwords
  changed at first sign-in, sign-out everywhere on a password change, and per-workshop appointments under
  separation-of-duties rules. All automated checks passed and it went live on 9 October; Android 0.0.16 was published.
\item \textbf{Requested accounts created}: the Assistant Director (Design) mailbox as Assistant Director (temporary
  password) and the Design Section mailbox as Ministry Admin (Google sign-in).
\item \textbf{Approval of finished reports}: the Ministry Admin can approve a workshop report, send it back or
  withdraw an approval with a reason, and hand it on to the Office; an approved report is locked.
\item \textbf{New features}: e-mail notices to designers when an inspector suggests a correction or sends a report back;
  sign-in with Microsoft and Yahoo accounts; inspectors and supervising officers can see a workshop's photographs
  and files; filters on the inspection list; reports can be previewed and downloaded without a connection;
  workshop join cards on the web; 3D models viewable on web and phone; inspection on the phone without signal.
\item \textbf{Plain wording}: every screen on web and phone was rewritten to tell officials what to do, without
  technical language.
\item \textbf{Platform}: the server was rebuilt on Ubuntu 26.04 with twice the memory and Python 3.14.8; Node.js 24 and
  the latest libraries throughout; database on PostgreSQL 17.11; all automated dependency updates applied.
\end{itemize}

\section{Next steps}
\begin{enumerate}[leftmargin=1.5em]
\item Onboard the Office's accounts, upload the annual plan, record sanction orders and begin fieldwork.
\item Make photographs and recordings private before real fieldwork; today anyone holding a file's link can open it.
\item Publish the next Android version with the 10 October features; e-mail delivery starts once Amazon approves
  the sending account (requested 10 October).
\item Owner decisions pending: whether a partial save of a workshop's first stage may clear details of a
  workshop opened from the annual plan, and whether the dashboard should name administrators holding
  workshop posts.
\end{enumerate}

\end{document}
